\documentclass[12pt]{article}
\usepackage[utf8]{inputenc}
\usepackage[T1]{fontenc}
\usepackage{hyperref}
\title{Maharastri Prakrit Language Resources for K-12 Curriculum Development}
\author{Compiled by Claude Sonnet 5.5 (claude-sonnet-5-5) for Andr\'as Kornai}
\date{}
\begin{document}
\maketitle

\section{Introduction}
Maharastri Prakrit (ISO 639-3 pmh; Glottocode maha1305) is a Middle Indo-Aryan literary language with no native speakers today. It is known from classical literature and from Jain writings. The search found no school curriculum, no primer and no revitalization programme for it. Any K-12 heritage curriculum would have to be built from the scholarly and reading materials listed below. All items here were seen in web search results.

\section{Primary Educational Resources}
No primer or textbook written for school use was found. The closest reading material is Jacobi's 1886 reader, which pairs Maharastri tales with a grammar and a dictionary. It is a scholarly reader, not a school textbook.
\begin{itemize}
\item Jacobi, \emph{Ausgew\"ahlte Erz\"ahlungen in M\=ah\=ar\=a\d{s}\d{t}r\=\i} (German text with transliteration, hosted on Jainpedia): \url{https://jainpedia.org/reading/ausgewaehlte-erzaehlungen-in-maharastri}
\end{itemize}

\section{Linguistic and Reference Materials}
\begin{itemize}
\item SOAS Digital Collections, Prakrit collection (grammar and dictionary items, including Hemacandra's grammar): \url{https://digital.soas.ac.uk/prakrit/all/table}
\item Wisdomlib Prakrit glossary, section ``m'' (general Prakrit, not specific to Maharastri): \url{https://www.wisdomlib.org/prakrit/glossary/m}
\item Jainpedia glossary entry on Maharastri Prakrit: \url{https://jainpedia.org/glossary/maha\d{r}a\d{s}\d{t}ri-prakrit}
\item Online Prakrit dictionary by ISJS, announced in a newsletter (an English version of the Prakrit--Hindi \emph{Paia-sadda-mahannavo}; general Prakrit, with uploading still in progress): \url{https://isjs-newsletter.in/Newsletter/2026/February/01/Prakrit.html}, \url{https://isjs-newsletter.in/Newsletter/2026/March/01/Prakrit.html}, \url{https://isjs-newsletter.in/Newsletter/2026/April/01/Prakrit.html}
\end{itemize}

\section{Religious and Community Resources}
Maharastri is the language of many Jain texts, so Jain scholarly bodies are the natural community partners. Jainpedia (above) is the only such resource found. No living community organisation dedicated to Maharastri was found.

\section{Digital and Audio Resources}
No audio, video or software for Maharastri was found in the searches. The digital resources are the scanned and online texts and glossaries listed above.

\section{Curriculum Development Context}
The language is extinct as a spoken language, so a curriculum would be a reading and heritage course: short tales, selected verses, basic grammar, and links to Jain literary history. The search did not find a Maharastri-specific grammar or dictionary online. Materials would need to be adapted from the general Prakrit works above. The language cannot be used for classroom conversation without new materials.
\end{document}